% ============================================================
% INTRODUCTION
% ============================================================

\section{Introduction}
\label{sec:introduction}

% HOOK: Why should the reader care? - Now with quantitative, cited facts
The cryptocurrency market has grown from a niche experiment to a multi-trillion-dollar asset class, with total market capitalization exceeding \$3.20 trillion as of December 2025 \cite{forbes2025crypto}, surpassing the previous peak of approximately \$3 trillion in November 2021 \cite{pessa2023age, watorek2023crypto}. This growth brings unique challenges for portfolio management: daily volatility that frequently exceeds traditional equity markets by several times \cite{bergsli2022volatility, naimy2021garch}, continuous 24/7 trading across fragmented exchanges with varying transaction costs \cite{barbon2021exchanges, kim2017transaction}, and an evolving asset universe as new tokens emerge and others fail. Furthermore, cross-asset correlations in cryptocurrency markets shift substantially between bull and bear regimes \cite{james2022correlations}, complicating covariance estimation for traditional optimization approaches.

% CONTEXT: Background on the problem
Deep reinforcement learning (DRL) has demonstrated strong performance in sequential decision-making domains from game playing to robotics, prompting researchers to explore its application to financial portfolio management \cite{jiang2017eIIE, lucarelli2020dqlcrypto}. The portfolio management problem---dynamically allocating capital across multiple assets to maximize risk-adjusted returns---naturally maps to a Markov Decision Process (MDP). At each decision point, the agent observes market conditions and its current holdings, then selects target portfolio weights subject to constraints such as position limits and transaction costs. Unlike supervised learning approaches that predict returns and then optimize, RL directly learns the allocation policy end-to-end, potentially capturing complex temporal dependencies and trading off immediate costs against future opportunities.

% GAP: What's missing in current approaches?
Despite promising results in controlled experiments, applying DRL to real-world portfolio management faces several challenges that prior work has incompletely addressed. First, most studies assume a fixed asset universe \cite{jiang2017eIIE, ye2020sarl}, but cryptocurrency portfolios must handle assets entering and leaving the tradable set as market conditions change. Second, the choice of action space---whether continuous weights, discrete allocation templates, or incremental adjustments---significantly affects learning dynamics but remains underexplored \cite{lucarelli2020dqlcrypto}. Third, rigorous comparison against strong traditional baselines is often lacking; the equal-weight (1/N) portfolio, despite its simplicity, has proven remarkably difficult to beat in out-of-sample tests \cite{demiguel2009optimal}.

% OBJECTIVE: What does this paper do?
This paper investigates whether deep reinforcement learning can learn profitable portfolio allocation strategies for cryptocurrency markets that outperform classical approaches. We develop a comprehensive experimental framework that addresses the practical challenges of variable asset universes, realistic transaction costs, and fair comparison across different algorithmic paradigms. Our goal is not merely to achieve positive returns---which depend heavily on the test period---but to understand under what conditions and architectural choices DRL methods succeed or fail relative to traditional baselines.

% APPROACH: How do we address it?
Our approach combines a realistic portfolio environment with multiple agent architectures. The environment handles monthly universe reconstitution, enforces trading constraints (turnover limits, concentration caps), and applies proportional transaction costs consistent with centralized exchange fee structures \cite{brauneis2022fees}. We implement four DRL agents spanning value-based and policy gradient paradigms: Deep Q-Network (DQN) with a novel delta-based action catalog, Double DQN (DDQN) to address overestimation bias \cite{vanHasselt2016ddqn}, REINFORCE for continuous policy optimization \cite{williams1992reinforce}, and REINFORCE with a learned value baseline to reduce gradient variance. These are benchmarked against three traditional strategies: equal weighting, market-capitalization weighting, and mean-variance optimization with shrinkage estimation \cite{ledoit2004honey}.

% CONTRIBUTIONS: What are the key contributions? - Consolidated from 5 to 4, merged reproducibility into first contribution
\paragraph{Contributions.} The main contributions of this work are:
\begin{enumerate}[leftmargin=*]
    \item \textbf{Variable universe environment}: A portfolio management environment handling dynamic asset membership with monthly reconstitution, cold-start rules, and automatic weight redistribution. We release a frozen dataset export enabling exact reproduction.
    
    \item \textbf{Delta-based discrete action space}: A 70-action catalog of relative portfolio adjustments for DQN agents, preserving portfolio continuity and enabling context-dependent feasibility learning.
    
    \item \textbf{Canonical padding for variable-size observations}: A scheme mapping variable-size observations to fixed-dimensional representations while preserving asset identity for Q-value learning.
    
    \item \textbf{Rigorous out-of-sample evaluation}: A 646-day held-out test demonstrating that DDQN outperforms passive baselines, and that policy gradient methods achieve competitive performance under deterministic evaluation.
\end{enumerate}

% OUTLINE: Paper structure
\paragraph{Paper organization.} The remainder of this paper is organized as follows. \Cref{sec:related_work} reviews related work in deep reinforcement learning for portfolio management and positions our contributions. \Cref{sec:methods} describes our environment design, agent architectures, and baseline implementations in detail. \Cref{sec:experiments} presents the experimental setup, training procedures, and quantitative results. \Cref{sec:discussion} interprets the divergent outcomes between value-based and policy gradient methods and discusses limitations. \Cref{sec:conclusion} concludes with directions for future research.
